\section{The calculated cycles in harmonic geometry}\label{sec:metrics}

The wall correspondences identify integral classes. Their metric
interpretation uses the harmonic metrics of the same compact
bundle moduli, with the normalization in
\cite{wuebben17,wuebben18}. We first calculate exact periods,
then state precisely which existing degeneration results apply.

\subsection{Periods, volumes and specialization}
Let $\omega_\eta$ be the smoothly varying Ricci-flat representative
of $H+sD_1+tD_2$. Fix the holomorphic determinant and use its
smoothly varying Hermitian--Einstein connections, together with the
real Dolbeault harmonic $L^2$ metric, unnormalized by the volume of $X$.
Its K\"ahler class is
\begin{equation}\label{eq:harmonic-class}
 [\Omega_\eta]=2\pi^2\theta_v(b_\eta),\qquad
 b_\eta=(0,[\omega_\eta],c_1(v)\cdot[\omega_\eta]/3),
\end{equation}
by \cite[\S2.4]{wuebben17}. Hence the period of a curve
functional $c$ is $2\pi^2(c,b_\eta)$ and a surface volume is
half its pairing with $[\Omega_\eta]^2$.

In chamber $A$ put $a=\zeta_0>0$, $b=\zeta_2>0$. Equations
\eqref{eq:levels}, \eqref{eq:curves} and
\eqref{eq:harmonic-class} give
\begin{equation}\label{eq:periods}
 \begin{gathered}
 \Area(c_0)=\pi a,\qquad \Area(c_2)=\pi b,\qquad
 \Area(f)=\pi(a+b),\\
 \Vol(P_0)=\tfrac12\pi^2a^2,\qquad
 \Vol(P_2)=\tfrac12\pi^2b^2.
 \end{gathered}
\end{equation}
The equality $f=c_0+c_2$ is the actual proper reduced
specialization from \cite[\S5.1]{wuebben17}. It is
preserved, with the signs of Table~\ref{tab:cycles}, by
the compact wall maps. In a ruled row the plane line and
common section have area $\pi|\zeta_k|$, and the ruling
has area $\pi|\zeta_1|$. The class $p_j-p_k$ therefore has
volume
\begin{equation}\label{eq:ruled-volumes}
 \Vol(\Sigma)=
 \begin{cases}
 \tfrac12\pi^2(\zeta_0^2-\zeta_2^2),&\sigma=B,E,\\
 \tfrac12\pi^2(\zeta_2^2-\zeta_0^2),&\sigma=C,F.
 \end{cases}
\end{equation}
Both expressions are positive in their sectors. Alternatively,
restrict $[\Omega_\eta]$ using
\eqref{eq:ruled-restriction} and compute its square on
$\FF_1$; this gives the same values.

For a wall action $h=\id$ or $s_{a_1}$ the exact period
identity is
\begin{equation}\label{eq:period-action}
 2\pi^2(hc,b_{\eta'})=2\pi^2(c,h^{-1}b_{\eta'}).
\end{equation}
This is a pairing identity, not an assertion that the wall
map is a metric isometry. In particular $h^{-1}b_{\eta'}$
need not be a polarization covector in the chosen two-parameter
family.

\subsection{Hypotheses of the metric comparison}
The analytic results used here have the following hypotheses,
verified for this family in \cite[\S9.1]{wuebben18}.
They are the conditions (G1)--(G3) of that paper:
\begin{enumerate}
\item The finitely many smooth projective chamber spaces equal
their generic Gieseker moduli and have proper surjective
maps to the connected normal projective central target.
There is a connected dense stable open, a finite Whitney
refinement of factor types, and agreement of the projective
and fixed-wall gauge topologies. The maps are isomorphisms
over the stable open. Every chamber object is $H$-semistable;
the central map assigns its $H$-graded class, and its complex
morphism is constant within a chamber.
\item Every central semistable sheaf is a bundle and every
central ideal HYM object has zero lost charge. Varying-metric
surface compactness gives smooth gauge subsequences and
excludes nearby charge.
\item The full convergent doubled Ext-quiver charts retain
all tangent factors, loops and multiplicities. Each complex
moment-zero fibre contains a module-simple point, meaning
a simple representation of its quiver with relations.
The classifying germs commute with the projective-to-affine
quotient and the central graded-class maps.
\end{enumerate}
The norm on every Ext space is its actual central harmonic
norm. At a central bundle $E_b=\bigoplus_iF_i\otimes\C^{m_i}$,
let $W_b=\Ext^1(E_b,E_b)$ be the full harmonic representation
and $K_b=(\prod_i\mathrm U(m_i))/\mathrm U(1)$ its effective
compact stabilizer. Denote the full quotient and its intrinsic
stratified length metric by
\[
 \cQ_{b,\nu}=\{\mu_{\C}=0,\ \mu_{\R}=\nu\}/K_b,
 \qquad d_{b,\nu}.
\]
Paths may be finitely piecewise smooth across representation-type
strata. Its radius $r_\nu$ is the norm of a balanced representative.
The moment levels use ordinary multiplicity trace:
\begin{equation}\label{eq:general-levels}
 \zeta_{b,i}(\eta)=2\pi\bigl(r_i\mu_\eta(v)-\deg_\eta F_i\bigr),
 \qquad r_i=\rank F_i,\qquad \sum_i m_i\zeta_{b,i}=0.
\end{equation}
For $L_1\oplus F$ the transverse two-arrow level is
$2\pi(2\mu_\eta(v)-\deg_\eta F)=-\zeta_1$.
Ordinary multiplicity trace is distinct from the rank-weighted
mass in small-spectrum statements.

For each fixed background the actual classifying chart
$\Phi_{b,\eta}$ satisfies the squared-metric estimate
\begin{equation}\label{eq:relative-metric}
 (1-\varepsilon_b)g_{b,\zeta_b(\eta)}
 \le\Phi_{b,\eta}^*g_\eta
 \le(1+\varepsilon_b)g_{b,\zeta_b(\eta)},\qquad
 \varepsilon_b=C_b(|\eta|+\|x\|^2).
\end{equation}
This is \cite[Corollary~2.2]{wuebben17}, with whole-preimage
coverage as in \cite[Proposition~5.1]{wuebben18}. The
constant depends on the fixed background and complementary
elliptic gap, but not on ratios of positive finite orbit
eigenvalues. A fixed deeper chart handles degeneration of
shallower backgrounds. No uniform gap is assumed for
separate backgrounds approaching a reducible bundle.

\subsection{Ambient collapse and complete rescalings}
The diameter statements in \cite[\S8]{wuebben17} use
ambient distances in the full compact moduli, not only
the intrinsic metrics of the exceptional surfaces. If
$\rho=a+b$ in chamber $A$, they give
\begin{equation}\label{eq:diameters}
 \diam_{g_\eta}P_i=\sqrt\rho\,(1+O(\rho))
                  D_i(a/\rho,b/\rho),
\end{equation}
where $D_i$ is the ambient diameter of the corresponding
plane in the complete normalized quotient. The finite
quotient diameters obey, uniformly in the positive ratio,
\[
 \sqrt{2a}\le\diam_{\mathrm{model}}P_0
          \le\frac{\pi}{\sqrt2}\sqrt a,
\]
and the same inequalities hold with $b$ for $P_2$.
The actual diameter has relative factor $1+O(\rho)$.
As $\eta\to0$, if $a/b\to0$, then $\diam P_0/\sqrt a\to\pi/2$;
the reciprocal statement holds for $P_2$.
For moving $\PP^1$ fibres over any compact subset of
$S^\circ$, $\diam/\sqrt{|\zeta_1|}\to\pi/2$ uniformly as
$\eta\to0$ through nonzero transverse levels $\zeta_1$.

The passage to ambient distances is an analytic input.
With squared-metric bounds $1/2,2$, a path starting at
radius $r_0<R$ costs at least $(R-r_0)/\sqrt2$ before
exiting radius $R$. This first-exit estimate confines
short competitors to the appropriate chart; comparison
of both length infima then yields the relative errors
even for a much smaller plane. These arguments apply
to our computed classes because their cycle realization
and classifying charts are the same ones.

We state the complete rescaling hypotheses to specify
the meaning of this compatibility. Fix a central chart
and suppress $b$. Suppose $\eta_n\to0$ is generic,
$\delta_n\downarrow0$, and
\begin{equation}\label{eq:normalized-level}
 \nu_n:=\zeta_b(\eta_n)/\delta_n^2\longrightarrow\nu.
\end{equation}
Choose balanced representatives $p_n$ at these normalized
levels which converge to $p$ after compact group adjustment.
Then \cite[\S8.3 and Theorem~8.1]{wuebben17} gives
\begin{equation}\label{eq:complete-limit}
 (M_{\sigma(\eta_n)},\delta_n^{-1}d_{g_{\eta_n}},
       \Phi_{b,\eta_n}([\delta_np_n]))
 \longrightarrow(\cQ_{b,\nu},d_{b,\nu},[p])
\end{equation}
in complete pointed Gromov--Hausdorff topology. On bounded
rescaled balls the relative error against the current
finite quotient is $O_L(|\eta_n|+\delta_n^2)$. This is
not a rate for $\cQ_{b,\nu_n}\to\cQ_{b,\nu}$.

For a further scale, let $x_*$ represent the limiting
point at level $\nu_* =\lim\nu_n$, with stabilizer $K$.
Decompose its full hyperk\"ahler normal representation as
$W=T\oplus W'$, where $T=W^K$ and $W'=T^\perp$ for the
retained invariant inner product. In this slice write
the current centers as $(u_n,z_n)\to0$ and require
\begin{equation}\label{eq:successive}
 \sigma_n\downarrow0,\qquad
 \frac{(\nu_n-\nu_*)|_{\Lie K}}{\sigma_n^2}\longrightarrow\lambda,
 \qquad z_n/\sigma_n\longrightarrow z_*.
\end{equation}
At physical scale $\delta_n\sigma_n$ the complete limit
is $T\times(W'\!\mathbin{/\!\!/\!/}_{\lambda}K)$,
based at $(0,[z_*])$. Tangential translation removes $u_n$;
no bound on $|u_n|/\sigma_n$ is required. The additional
slice error on a fixed bounded ball is
$O_{x_*}(|\nu_n-\nu_*|+|u_n|+\sigma_n R)$.
Inverse coverage and the outer and inner exit costs
give whole actual balls at each stage. This is the
finite-stage theorem of \cite[Theorem~8.1]{wuebben17}
and \cite[Proposition~8.3]{wuebben18}; any fixed finite
iteration must retain the same normalized-level and
normal-basepoint data at every stage.

At a nonzero $\zeta_0$ or $\zeta_2$ wall the normal
resolution is $T^*\PP^2$. At a nonzero $\zeta_1$ wall
the full model has a harmonic $\C^2$ tangent factor
and the $A_1$ transverse quotient, resolved by the
Eguchi--Hanson metric at nonzero level. An $\FF_1$
component belongs to the interacting quotient; it is
not asserted to have a complete cotangent-$\FF_1$
limit. Table~\ref{tab:cycles} identifies the effective
cycles in each of these inherited models.

\subsection{Global convergence with the same contractions}
Under the three stated hypotheses, $(Z,d_0)$ is the
compact completion of the central harmonic metric on
$U$. We recall the quantitative form of
\cite[Theorem~5.5]{wuebben18}, since it applies to the
same $\pi_\eta$ used in our integral calculation.
Choose $\varepsilon_j\downarrow0$, $\varepsilon_j<1/8$.
For each accuracy choose a finite central chart cover
with outer radii $16r_{ji}$, inner radii $r_{ji}/4$,
and squared-metric error at most $\varepsilon_j$.
Then shrink a parameter neighborhood $\mathcal P_j$
so that whole-preimage coverage holds.

For a finite quotient contraction $q_\zeta:\cQ_\zeta\to\cQ_0$
retain the distance discrepancy
\[
 b_{i,r}(\eta)=\sup_{\substack{r_0(z),r_0(z')\le r\\
                    q_\zeta(x)=z,\ q_\zeta(x')=z'}}
       |d_\zeta(x,x')-d_0(z,z')|
\]
and both radial discrepancies
\[
 \begin{aligned}
 a_{i,r}(\eta)&=\sup_{\substack{q_\zeta(x)=z\\r_0(z)\le r}}
                     |r_\zeta(x)-r_0(z)|,\\
 \widetilde a_{i,r}(\eta)&=\sup_{r_\zeta(x)\le r}
                     |r_0(q_\zeta(x))-r_\zeta(x)|.
 \end{aligned}
\]
The symbols $d_0,r_0$ in these formulas are the finite
zero-quotient distance and radius. Properness and
intrinsic parameter continuity make all three finite
and tending to zero. Arrange $a,\widetilde a<r/16$
and $b<r$. Choose
\[
 0<\ell_j<\min_i\frac{r_{ji}}{16\sqrt2},\quad
 \alpha_j=\sqrt{\frac{1+\varepsilon_j}{1-\varepsilon_j}}-1,
 \quad\beta_j(\eta)=\sqrt{1+\varepsilon_j}
                         \max_i b_{i,r_{ji}}(\eta).
\]
A fixed finite central cover supplies
$D_*\ge\diam(Z,d_0)$. The stated theorem then gives
\begin{equation}\label{eq:global-distortion}
 \dis(\operatorname{graph}\pi_\eta)
 \le\inf_{j:\eta\in\mathcal P_j}
 \left[\left(\alpha_j+\frac{\beta_j(\eta)}{\ell_j}\right)D_*
                  +\beta_j(\eta)\right]\longrightarrow0.
\end{equation}
The Gromov--Hausdorff distance is bounded by half this
quantity. The proof uses the two radial discrepancies
to compare short ambient pairs in both directions,
and then subdivides paths into lengths at most $\ell_j$.
For fixed $j$, first let $\eta\to0$; then let the
accuracy index tend to infinity. No universal power
law is asserted for this modulus.

The period calculation and the cycle table identify
the integral classes entering these metric results.
They leave the harmonic norms, no-charge compactness,
convergent charts and normalized basepoint data as
analytic hypotheses. Bubbling, unrestricted HYM
dynamics and genuine complex smoothings require
additional results.
